\documentclass[11pt]{article}
\usepackage[margin=1in]{geometry}
\usepackage{amsmath,amssymb,amsthm}
\usepackage{graphicx}
\usepackage[hidelinks]{hyperref}
\setlength{\emergencystretch}{2em}

\newtheorem{theorem}{Theorem}
\newtheorem{lemma}{Lemma}
\newtheorem{remark}{Remark}

\newcommand{\C}{\mathbb C}
\newcommand{\Kcal}{\mathcal K}
\newcommand{\dist}{\operatorname{dist}}
\newcommand{\diag}{\Delta}

\title{The Sequence-Quotient Reflection Theorem Fails for Compact Operators}
\author{Run \texttt{fa\_banach\_001}, lane 19}
\date{June 17, 2026}

\begin{document}
\maketitle

\section*{Source Problem and Internal Lemma}

The source paper is Jinghao Huang, Karimbergen Kudaybergenov, and Rui Liu,
\emph{The ball-covering property of von Neumann algebras and noncommutative
symmetric spaces}, arXiv:2311.14257.  After proving the commutative and von
Neumann boundary cases, it asks for a noncommutative C*-algebra
classification:

\begin{center}
\includegraphics[width=0.92\textwidth]{figures/open_problem_crop.png}
\end{center}

This packet is not a counterexample to the source problem.  It is a
counterexample to an intermediate reflection lemma proposed during the
lane-19 push:

\begin{quote}
If a countable BCP witness \(C\subset A\) is captured by a unit vector in
\(\ell_\infty(A)/c_0(A)\), then \(C\) is captured by a unit vector in \(A\).
\end{quote}

\section*{Statement}

Let \(H=\ell_2\), let \(A=\Kcal(H)\), and let
\[
        A_\infty=\ell_\infty(A)/c_0(A).
\]
Let \(\diag:A\to A_\infty\) denote the diagonal embedding.

\begin{theorem}
\label{thm:main}
There is a countable BCP witness \(C\subset S_A\) and an element
\(v\in S_{A_\infty}\) such that
\[
        \dist(\diag(c),\C v)=\|c\|
        \qquad(c\in C),
\]
but there is no \(u\in S_A\) such that
\[
        \dist(c,\C u)=\|c\|
        \qquad(c\in C).
\]
\end{theorem}

\section*{Idea of the Proof}

Let \(p_n\) be the rank-one projection onto the \(n\)-th basis vector of
\(\ell_2\).  Every compact operator is asymptotically disjoint from these
rank-one projections:
\[
        \|a-\lambda p_n\|\longrightarrow \max\{\|a\|,|\lambda|\}.
\]
Thus \(v=[(p_n)]\) captures the diagonal copy of every compact operator in
the sequence quotient.

On the other hand, if \(C\) is dense in the unit sphere of \(A\), then no
actual unit compact operator \(u\) can capture all of \(C\), because some
\(c\in C\) is close to \(u\), giving
\[
        \dist(c,\C u)<\|c\|.
\]

\section*{The Compact-Operator Limit}

\begin{lemma}
\label{lem:limit}
Let \((e_n)\) be the standard orthonormal basis of \(\ell_2\), and let
\(p_n\) be the rank-one projection onto \(\C e_n\).  For every
\(a\in\Kcal(\ell_2)\) and every \(\lambda\in\C\),
\[
        \lim_{n\to\infty}\|a-\lambda p_n\|
        =
        \max\{\|a\|,|\lambda|\}.
\]
\end{lemma}

\begin{proof}
Let \(P_N\) be the projection onto
\(\operatorname{span}\{e_1,\ldots,e_N\}\).  Since \(a\) is compact,
\[
        \|a-P_NaP_N\|\longrightarrow0.
\]
Fix \(\varepsilon>0\), and choose \(N\) with
\[
        \|a-P_NaP_N\|<\varepsilon .
\]
For \(n>N\), the operators \(P_NaP_N\) and \(p_n\) act on orthogonal summands,
so
\[
        \|P_NaP_N-\lambda p_n\|
        =
        \max\{\|P_NaP_N\|,|\lambda|\}.
\]
Therefore, for \(n>N\),
\[
\begin{aligned}
 \bigl|\|a-\lambda p_n\|
       -\max\{\|P_NaP_N\|,|\lambda|\}\bigr|
 &\le \|a-P_NaP_N\|\\
 &<\varepsilon .
\end{aligned}
\]
Also \(|\|P_NaP_N\|-\|a\||<\varepsilon\) after increasing \(N\), since
\(P_NaP_N\to a\) in norm.  Letting \(\varepsilon\downarrow0\) proves the
claim.
\end{proof}

\section*{Proof of the Counterexample}

\begin{proof}[Proof of Theorem~\ref{thm:main}]
Let \(C\subset S_A\) be a countable dense subset of the unit sphere of
\(A=\Kcal(\ell_2)\).  This is a BCP witness: for every \(u\in S_A\), choose
\(c\in C\) with \(\|c-u\|<1/2\).  Then
\[
        \dist(c,\C u)\le \|c-u\|<1/2<1=\|c\|.
\]

Let \(v\in A_\infty\) be the class of the bounded sequence \((p_n)\).  Since
\(\|p_n\|=1\), we have \(v\in S_{A_\infty}\).  For any \(a\in A\) and
\(\lambda\in\C\), Lemma~\ref{lem:limit} gives
\[
        \|\diag(a)-\lambda v\|
        =
        \limsup_n\|a-\lambda p_n\|
        =
        \max\{\|a\|,|\lambda|\}.
\]
Taking the infimum over \(\lambda\in\C\), we obtain
\[
        \dist(\diag(a),\C v)=\|a\|.
\]
In particular this holds for every \(c\in C\).

It remains to show that no actual unit vector in \(A\) reflects this capture.
Let \(u\in S_A\).  Since \(C\) is dense in \(S_A\), choose \(c\in C\) with
\[
        \|c-u\|<1/2.
\]
Then
\[
        \dist(c,\C u)\le \|c-u\|<1/2<1=\|c\|.
\]
Thus \(u\) does not capture \(C\).  Since \(u\in S_A\) was arbitrary, no such
reflecting unit vector exists in \(A\).
\end{proof}

\section*{Consequence for the C*-BCP Program}

This counterexample shows that diagonal capture in
\(\ell_\infty(A)/c_0(A)\) is too weak to imply actual capture in \(A\), even
for a countable BCP witness.  Therefore the proposed reflection shortcut
cannot prove
\[
        \text{C*-algebra BCP}\Rightarrow\text{C*-algebra UBCP}.
\]

There is no contradiction with known positive results: \(A=\Kcal(\ell_2)\)
is separable, hence has UBCP.  The point is subtler.  UBCP spaces can still
have sequence-quotient elements that capture countable diagonal sets.
Therefore a proof of BCP\(\Rightarrow\)UBCP must use more than the mere
existence of sequence-quotient capture obtained by diagonalization.

\section*{Novelty and Search Bounds}

Local indexes were searched for arXiv:2311.14257, reflection theorem,
sequence quotient, countable capture, compact operators, BCP, and UBCP.  The
previous lane-19 attempts proposed the reflection route but did not record
this compact-operator counterexample.  This packet is a counterexample to an
agent-generated intermediate lemma, not a new answer to the source paper's
Question 3.

\section*{Human Review Focus}

Reviewers should check:
\begin{enumerate}
    \item the compact approximation step
          \(P_NaP_N\to a\) for \(a\in\Kcal(\ell_2)\);
    \item the orthogonal-block norm formula for
          \(P_NaP_N-\lambda p_n\) when \(n>N\);
    \item the computation of distance to the line \(\C v\) in the quotient;
    \item the fact that a dense subset of the unit sphere is a BCP witness
          but cannot be actually captured by any unit vector.
\end{enumerate}

\begin{thebibliography}{9}

\bibitem{HKL}
Jinghao Huang, Karimbergen Kudaybergenov, and Rui Liu,
\emph{The ball-covering property of von Neumann algebras and noncommutative
symmetric spaces},
arXiv:2311.14257, 2023.

\bibitem{Conway}
John B. Conway,
\emph{A Course in Functional Analysis},
Springer, 1990.

\bibitem{lane19capture}
Run \texttt{fa\_banach\_001}, lane 19,
\emph{BCP and Countable Capture Are Exact Opposites},
solution packet
\texttt{2311.14257\_bcp\_countable\_capture\_exact\_equivalence}, 2026.

\bibitem{lane19reflection}
Run \texttt{fa\_banach\_001}, lane 19,
\emph{Local Transition Bumps and the Reflection Wall},
attempt note
\texttt{2311.14257\_local\_transition\_bump\_reflection\_push}, 2026.

\end{thebibliography}

\end{document}

